\section{Signal Processing Algorithms}

\subsection{Channelization via FFT}

The F-Engine applies the Discrete Fourier Transform:

\begin{equation}
    X[k] = \sum_{n=0}^{N-1} x[n] \, e^{-2\pi j kn / N}
\end{equation}

NumPy's FFT computes this in $\mathcal{O}(N \log N)$ operations. Window functions reduce spectral leakage: Hanning ($-31$ dB sidelobes), Hamming ($-43$ dB), Blackman ($-58$ dB), or rectangular (no tapering). Window selection trades resolution against leakage suppression.

Overlapping windows reduce scalloping loss and improve time resolution at computational cost. Channelized output maps bins to frequencies: $f[k] = k \cdot f_s / N$.

Standard FFT ordering places the zero-frequency (DC) component at $k=0$, positive frequencies at $k=1$ to $k=N/2-1$, the Nyquist frequency at $k=N/2$, and negative frequencies at $k=N/2+1$ to $k=N-1$. Imaging applications often rearrange this to place DC in the centre using \texttt{numpy.fft.fftshift}, but the correlator outputs frequencies in native FFT order for efficiency.

\subsection{Geometric Delay Compensation}

Signals from a distant source reach different antennas at different times due to path-length differences. For a source at direction $\hat{s}$ (unit vector), an antenna at position $\vec{r}_i$ relative to the array centre experiences a geometric delay:

\begin{equation}
    \tau_i = \frac{\vec{r}_i \cdot \hat{s}}{c}
\end{equation}

where $c$ is the speed of light. This delay varies between antennas; the correlator compensates by applying a phase correction that aligns all antennas to a common reference.

The delay compensation operates in the frequency domain. A time delay $\tau$ corresponds to a phase shift $\phi = -2\pi f \tau$ at frequency $f$. The corrected spectrum is:

\begin{equation}
    \tilde{X}_i[k] = X_i[k] \cdot e^{-2\pi j f[k] \tau_i}
\end{equation}

The implementation pre-computes the phase angles for all antennas and channels, then applies them via element-wise multiplication of spectral arrays with a complex exponential array. This operation scales as $\mathcal{O}(MAK)$ where $M$ is the number of antennas, $A$ is the number of spectra in a chunk, and $K$ is the number of channels—a modest computational burden compared to FFT or correlation.

Antenna positions are specified in a local coordinate frame (typically East-North-Up or X-Y-Z Cartesian). The implementation supports both 2D arrays (antennas in a horizontal plane) and 3D configurations. For 2D arrays, the third coordinate is implicitly zero. The phase centre $\hat{s}$ is provided as a unit vector; for a source at altitude $\theta$ and azimuth $\phi$, this converts to:

\begin{equation}
    \hat{s} = \begin{pmatrix} \cos\theta \cos\phi \\ \cos\theta \sin\phi \\ \sin\theta \end{pmatrix}
\end{equation}

In practice, the implementation accepts either Cartesian unit vectors or will convert from spherical coordinates if provided.

The delay model assumes a single fixed phase centre. Extensions for tracking sources or multiple phase centres would involve updating $\tau_i$ between chunks based on time-dependent source positions, but the current implementation uses static delays suitable for snapshot observations.

\subsection{Cross-Correlation and Visibility Computation}

The X-Engine computes cross-correlation products, known as visibilities, between every pair of antennas. For antennas $i$ and $j$ at frequency channel $k$, the visibility is:

\begin{equation}
    V_{ij}[k] = X_i[k] \cdot X_j^*[k]
\end{equation}

where $^*$ denotes complex conjugation. Autocorrelations (where $i=j$) yield:

\begin{equation}
    V_{ii}[k] = |X_i[k]|^2
\end{equation}

which is real and represents the power spectrum of antenna $i$.

The baseline ordering scheme places autocorrelations first, followed by cross-correlations in lexicographic order. For $M$ antennas:

\begin{enumerate}
    \item Baselines 0 to $M-1$: autocorrelations $(0,0), (1,1), \ldots, (M-1,M-1)$
    \item Baselines $M$ onwards: cross-correlations $(0,1), (0,2), \ldots, (M-2,M-1)$
\end{enumerate}

This produces $M + M(M-1)/2 = M(M+1)/2$ baselines total. For a four-antenna system, this yields 4 autocorrelations and 6 cross-correlations, totalling 10 baselines.

The implementation computes visibilities using nested loops over baselines and channels. While this could be vectorised further using outer product operations, the double loop proves adequate for the modest number of baselines. Each visibility computation involves one complex multiplication (or magnitude-squared for autocorrelations), requiring six floating-point operations per baseline per channel.

Time integration averages visibilities over multiple spectra to reduce thermal noise. The correlator accumulates visibilities in buffers:

\begin{equation}
    \text{Acc}_{ij}[k] = \sum_{t=0}^{T-1} V_{ij}^{(t)}[k]
\end{equation}

where $T$ is the number of spectra spanning the integration time. After accumulating $T$ spectra, the average is:

\begin{equation}
    \langle V_{ij}[k] \rangle = \frac{1}{T} \text{Acc}_{ij}[k]
\end{equation}

The number of spectra $T$ is determined by the configured integration time $t_{\text{int}}$ and the time per spectrum $\Delta t_{\text{spec}} = N / f_s$:

\begin{equation}
    T = \left\lfloor \frac{t_{\text{int}}}{\Delta t_{\text{spec}}} \right\rfloor = \left\lfloor \frac{t_{\text{int}} f_s}{N} \right\rfloor
\end{equation}

For example, with $N=256$ channels, $f_s=1024$ Hz, and $t_{\text{int}}=1$ second, one integrates $T = \lfloor 1 \times 1024 / 256 \rfloor = 4$ spectra.

Autocorrelations naturally produce real values, but numerical errors in complex arithmetic can introduce small imaginary components. The implementation explicitly takes the real part of autocorrelations after averaging to ensure the output is strictly real.

\subsection{Computational Complexity Analysis}

The overall computational cost per data chunk breaks down by stage. For a chunk of $C$ samples across $M$ antennas, channelized into $K$-point FFTs with stride $S$, the number of spectra produced is $A = \lfloor (C - K) / S \rfloor + 1$.

\textbf{F-Engine}: Each antenna requires $A$ FFTs of $K$ points. The FFT cost is $\mathcal{O}(K \log K)$ per transform, so total cost is $\mathcal{O}(MAK \log K)$. For $M=4$, $A=16$, $K=256$: approximately $4 \times 16 \times 256 \times \log_2(256) = 4 \times 16 \times 256 \times 8 \approx 131{,}000$ operations.

\textbf{Delay Engine}: Applying phase corrections requires one complex multiply per antenna per spectrum per channel: $\mathcal{O}(MAK)$. For the same parameters: $4 \times 16 \times 256 \approx 16{,}400$ complex multiplies (about 100{,}000 floating-point operations).

\textbf{X-Engine}: Cross-correlation involves $B = M(M+1)/2$ baselines, each requiring one complex multiply per channel per spectrum: $\mathcal{O}(BAK)$. For $M=4$ ($B=10$), $A=16$, $K=256$: $10 \times 16 \times 256 \approx 41{,}000$ complex multiplies (about 250{,}000 operations).

The F-Engine dominates computational cost due to the FFT. The delay and correlation stages add relatively modest overhead. This cost structure confirms that the FX architecture is appropriate: for high spectral resolution (large $K$), the $K \log K$ FFT cost grows slower than the $K^2$ cost of time-domain convolution in XF architectures.

Memory requirements scale with array dimensions. Storing one chunk's time-domain data requires $16 \cdot M \cdot C$ bytes (assuming 8-byte complex samples). Channelized data requires $16 \cdot M \cdot A \cdot K$ bytes. Accumulated visibilities need $16 \cdot B \cdot K$ bytes. For $M=4$, $C=4096$, $K=256$, $A=16$, $B=10$: time data is 256 KB, channelized data is 256 KB, and visibilities are 40 KB. Total working set remains well under 1 MB per chunk, easily fitting in modern CPU caches for efficient processing.
